\documentclass{Math_Note}
\setlength{\parindent}{2em}

\title{\textbf{Mathmatical Analysis Zorich Comprehensions}}
\author{Buce-Ithon}
\newdateformat{mydate}{\twodigit{\THEDAY}{ }\shortmonthname[\THEMONTH], \THEYEAR}
\date{\today}

\begin{document}
% Title page
\maketitle

% Content
\newpage
\tableofcontents
\newpage

% Surface-figure
\begin{figure}[H]
    \centering
    
\includegraphics[scale=0.24]{"./Figures/preface.jpg"}
    \caption{One week, one Zorich~}
\end{figure}

% \setcounter{section}{-1}
\newpage
\section{Multiple Integrals Short Sheet (Calculus)}
\subsection{Development of $1$-dim Reimann Integral-Line Integral}
\subsubsection{Intrduction ($ds$)}
\fbox{\begin{minipage}{0.95\textwidth}
\textcolor{skycyan}{Introduction}\\
$I=\int_{C} f(x,y) \, ds,\ C:r=r(\theta),\ \theta \in [\theta_{1},\theta_{2}]$\\
$C:\left\{\begin{matrix} x=r(\theta)\cos\theta \\ y=r(\theta)\sin\theta \end{matrix}\right. ,\ f(x,y)=f(r\cos\theta,r\sin\theta)$\\
$ds=\sqrt{ x'(\theta)^{2}+y'(\theta)^{2} }d\theta=\sqrt{ r'^{2}+r^{2} }d\theta$\\
$I=\int_{\theta_{1}}^{\theta_{2}} f(r\cos\theta,r\sin\theta)\sqrt{ r'^{2}+r^{2} } \, d\theta$
\end{minipage}}
\subsubsection{I-Line Integral}
\fbox{\begin{minipage}{0.95\textwidth}
\textcolor{skycyan}{Definition}\\
$I=\int_{C} f(x,y) \, ds=\left[\left\{\begin{matrix} x=\varphi(t) \\ y=\psi(t) \end{matrix}\right.\right]=\int_{D_{t}} f(\varphi(t),\psi(t))\sqrt{ \varphi'(t)^{2}+\psi'(t)^{2} } \, dt$
\end{minipage}}
\newline\newline
\fbox{\begin{minipage}{0.95\textwidth}
\textcolor{skycyan}{A mean value property}\\
$\rho(x,y)=f(x,y),\ \int_{C(A,B)} f(x,y) \, ds=M,\ \int_{C(C,D)} f(x,y) \, ds=m$\\
$\frac{\int_{C(C,D)} f(x,y) \, ds}{\int_{C(A,B)} f(x,y) \, ds} = \frac{\int_{C(C,D)} C \, ds}{\int_{C(A,B)} C \, ds} = \frac{S_{CD}}{S_{AB}}$
\begin{figure}[H]
    \centering
    \includegraphics[scale=0.36]{"./Figures/250809_1.png"}
\end{figure}
\end{minipage}}
\subsubsection{II-Line Integral}
\fbox{\begin{minipage}{0.95\textwidth}
\textcolor{skycyan}{Definition}\\
$I=\int_{C} f(x,y) \, dx=[x=\varphi(t),y=\psi(t)]=\int_{D_{t}} f(\varphi(t),\psi(t))\varphi'(t) \, dt$\\
$I=\int_{C} f(x,y) \, dy=[x=\varphi(t),y=\psi(t)]=\int_{D_{t}} f(\varphi(t),\psi(t))\psi'(t) \, dt$\\
$I=\int_{C} P \, dx+Q \, dy=[x=\varphi(t),y=\psi(t)]=\int_{D_{t}} [P\varphi'+Q\psi'] \, dt$\\
\end{minipage}}
\subsubsection{Calculate area using II-Line Integral}
\fbox{\begin{minipage}{0.95\textwidth}
\textcolor{skycyan}{Calculate}\\
$D=-\int_{C} y \, dx=\int_{C} x \, dy=\frac{1}{2}\int_{C} x \, dy - y \, dx$, where $C$ is a simple closed cure with positive direction.
\begin{figure}[H]
    \centering
    \includegraphics[scale=0.36]{"./Figures/250809_2.png"}
\end{figure}
\end{minipage}}
\subsubsection{Conditions for curve integrals to be independent of the path}
\fbox{\begin{minipage}{0.95\textwidth}
\textcolor{skycyan}{Theorem}\\
In simply connected space $D$ and given a path $C:A\to B$, II-line integral $\int_{C} P \, dx+Q \, dy$ is dependent of the integration path $C:A\to B$ 
i.f.f. $\frac{\partial P}{\partial y}=\frac{\partial Q}{\partial x}$.
\end{minipage}}
\newline\newline
\fbox{\begin{minipage}{0.95\textwidth}
\textcolor{skycyan}{Proof idea}\\
Necessity "$\Rightarrow$": \\
By studying the differentiation method (partial differentiation with respect to $x$ and $y$) of $F(x, y) = \int_{(x_0, y_0)}^{(x, y)} P  dx + Q  dy$ and the primitive function of $P  dx + Q  dy$ (regarding it as the total differential of some function) (i.e., $\exists \Phi$ such that $P = \frac{\partial \Phi}{\partial x}$, $Q = \frac{\partial \Phi}{\partial y}$ and $\int_C P  dx + Q  dy = \Phi(x, y) \big|_{(x_A, y_A)}^{(x_B, y_B)}$), 
we can conclude that the necessary and sufficient condition for $\int_C P  dx + Q  dy$ to be path-independent is that $P  dx + Q  dy$ ($= d\Phi(x)$) is an exact differential. The latter implies $\frac{\partial P}{\partial y} = \frac{\partial Q}{\partial x}$ (necessity). \\
Sufficiency "$\Leftarrow$": \\
Fix $y \in [c, d]$, integrate $P = \frac{\partial \Phi}{\partial x}$ on $[x_0, x] \subset D$, then $\Phi(x, y) = \int_{x_0}^x P  dx + \Phi(x_0, y)$.
Also, fix $x = x_0$, integrate $Q = \frac{\partial \Phi}{\partial y}$ on $[y_0, y] \subset D$, then $\Phi(x_0, y) = \int_{y_0}^y Q(x_0, y)  dy + \Phi(x_0, y_0)$.\\
Then $\Phi(x, y) = \int_{x_0}^x P(x, y)  dx + \int_{y_0}^y Q(x_0, y)  dy + \Phi(x_0, y_0)$ or $\Phi(x, y) = \int_{y_0}^y Q(x, y)  dy + \int_{x_0}^x P(x, y_0)  dx + \Phi(x_0, y_0)$, which satisfies $\frac{\partial P}{\partial y} = \frac{\partial Q}{\partial x}$.
\end{minipage}}
\subsection{$2$-dim Multiple Integral}
\subsubsection{Definition}
\fbox{\begin{minipage}{0.95\textwidth}
\textcolor{skycyan}{Definition}\\
$\iint_{P} f(x,y) dP=\lim\limits_{ \lambda=\max_{i}\{P_{i}\} \to 0 } \sum\limits_{i} f(\xi_{i},\eta_{i})\Delta P_{i}=V$
\begin{figure}[H]
    \centering
    \includegraphics[scale=0.32]{"./Figures/250809_3.png"}
\end{figure}
\end{minipage}}
\subsubsection{Calculate: $2$-dim Multiple Integral $\to$ Iterated Integral}
\fbox{\begin{minipage}{0.95\textwidth}
\textcolor{skycyan}{$2$-dim Multiple Integral $\to$ Iterated Integral}\\
\underline{Rectangle $P=[a,b]\times[c,d]$}: \\
If $\iint_P f(x,y) dP$ exists and $I_1=\int_a^b f(x,y) dx$ exists $\forall y \in [c,d]$, then
$\iint_P f(x,y) dP=\int_c^d dy \int_a^b f(x,y) dx\ .$\\
If $\iint_P f(x,y) dP$ exists and $I_2=\int_c^d f(x,y) dy$ exists $\forall x \in [a,b]$, then
$\iint_P f(x,y) dP=\int_a^b dx \int_c^d f(x,y) dy\ .$ \\
$\implies$ If $I_1$ and $I_2$ exist, then $\iint_P f(x,y) dP=\int_a^b dx \int_c^d f(x,y) dy=\int_c^d dy \int_a^b f(x,y) dx$.\\
\underline{Curved trapezoid $P=[a,b]\times[y_0(x),Y(x)]$ or $P=[x_0(y),X(y)]\times[c,d]$}: \\
If $\iint_P f(x,y) dP$ exists and $I_2=\int_{y_0(x)}^{Y(x)} f(x,y) dy$ exists $\forall x \in [a,b]$, then
$\iint_P f(x,y) dP=\int_a^b dx \int_{y_0(x)}^{Y(x)} f(x,y) dy\ .$\\
If $\iint_P f(x,y) dP$ exists and $I_1=\int_{x_0(y)}^{X(y)} f(x,y) dx$ exists $\forall y \in [c,d]$, then
$\iint_P f(x,y) dP=\int_c^d dy \int_{x_0(y)}^{X(y)} f(x,y) dx\ .$ \\
$\implies$ If $I_1$ and $I_2$ exist, then $\iint_P f(x,y) dP=\int_a^b dx \int_{y_0(x)}^{Y(x)} f(x,y) dy=\int_{y_0(x)}^{Y(x)} dy \int_a^b f(x,y) dx$.
\begin{figure}[H]
    \centering
    \includegraphics[scale=0.32]{"./Figures/250815_1.png"}
\end{figure}
\end{minipage}}
\subsubsection{Green's Formula: $2$-dim Multiple Integral $\longleftrightarrow$ II-Line Integral}
\fbox{\begin{minipage}{0.95\textwidth}
\textcolor{skycyan}{Green's Formula}\\
$$\iint_{D} (\frac{\partial P}{\partial x}-\frac{\partial Q}{\partial y}) \, dx dy=\oint_{C=\partial D} P \, dy+Q \, dx$$
or 
$$\iint_{D} (\frac{\partial P}{\partial x}+\frac{\partial Q}{\partial y}) \, dx dy=\oint_{C=\partial D} P \, dy-Q \, dx$$
where $D$ is a simply connected domain in $\mathbb{R}^2$ with boundary $C=\partial D$ and $P, Q$ are continuously differentiable functions on $D$.
\end{minipage}}
\newline\newline
\fbox{\begin{minipage}{0.95\textwidth}
\textcolor{skycyan}{Proof idea}\\
$P, Q, \frac{\partial P}{\partial x}, \frac{\partial Q}{\partial y}$ continuously on $D$, $\partial D \triangleq C$.\\
$\iint_{D} \frac{\partial P}{\partial x} dx dy = \int_c^d dy \int_{x_0(y)}^{X(y)} \frac{\partial P}{\partial x} dx = \int_c^d [P(X(y), y) - P(x_0(y), y)] dy = \oint_{\partial D} P(x,y) dy$\\
$\iint_{D} \frac{\partial Q}{\partial y} dx dy = \int_a^b dx \int_{y_0(x)}^{Y(x)} \frac{\partial Q}{\partial y} dy = \int_a^b [Q(x, Y(x)) - Q(x, y_0(x))] dx = - \oint_{\partial D} Q(x,y) dx$
\begin{figure}[H]
    \centering
    \includegraphics[scale=0.36]{"./Figures/250815_2.png"}
\end{figure}
\end{minipage}}
\subsubsection{Variable Substitution in Double Integrals}
\fbox{\begin{minipage}{0.95\textwidth}
\textcolor{skycyan}{Intro}\\
$D=\oint_{S=\partial D} x dy=\left[\begin{cases} x=x(\xi,\eta)\\ y=y(\xi,\eta) \end{cases}\right]=\oint_{\Sigma=\partial\Delta} x(\xi,\eta)\frac{\partial y}{\partial \xi} d\xi + x(\xi,\eta)\frac{\partial y}{\partial \eta} d\eta=\iint_{\Delta} \frac{D(x,y)}{D(\xi,\eta)} d\xi d\eta$
\end{minipage}}
\newline\newline
\fbox{\begin{minipage}{0.95\textwidth}
\textcolor{skycyan}{Analysis}\\
As we could see, the essential of $\begin{cases} x=x(\xi,\eta)\\ y=y(\xi,\eta) \end{cases}$ actually is coordinate transform between the $x,y$-plane and the $\xi,\eta$-plane.\\
$D=\oint_S x dy=\left[\begin{cases} x=x(\xi,\eta)\\ y=y(\xi,\eta) \end{cases},\begin{cases} \xi=\xi(t)\\ \eta=\eta(t) \end{cases} \right]=\int_{\alpha}^{\beta}x(\xi(t),\eta(t))(\frac{\partial y}{\partial\xi}\xi'(t)+\frac{\partial y}{\partial\eta}\eta'(t))dt=\oint_{\Sigma}x(\xi,\eta)\frac{\partial y}{\partial\xi}d\xi+x(\xi,\eta)\frac{\partial y}{\partial\eta}d\eta$,
let $x(\xi,\eta)\frac{\partial y}{\partial\xi} \triangleq P(x,y),\ x(\xi,\eta)\frac{\partial y}{\partial\eta} \triangleq Q(x,y)$, then $D=\oint_{\Sigma} P d\xi + Q d\eta=(\frac{\partial P}{\partial\eta}=\frac{\partial Q}{\partial\xi},\ \text{using Greeni's formula})= \iint_{\Delta}(\frac{\partial Q}{\partial \xi}-\frac{\partial P}{\partial \eta}) d\xi d\eta=\iint_{\Delta}\frac{D(x,y)}{D(\xi,\eta)} d\xi d\eta$.
\end{minipage}}
\newline\newline
\fbox{\begin{minipage}{0.95\textwidth}
\textcolor{skycyan}{Figures for aforementioned analysis}\\
\begin{figure}[H]
    \centering
    \includegraphics[scale=0.32]{"./Figures/250815_3.png"}
\end{figure}
\end{minipage}}
\newline\newline
\fbox{\begin{minipage}{0.95\textwidth}
\textcolor{skycyan}{Idea of Analysis}: \\
Using $\begin{cases} x=x(\xi,\eta)=x(\xi(t),\eta(t))\\ y=y(\xi,\eta)=y(\xi(t),\eta(t)) \end{cases}$ (coordinate transform with parameterization), transform the II-line integration $\oint_Sx dy$ (or $-\oint_S y dx$, or $\frac{1}{2}\oint_S x dy - y dx$) on $x,y$-plane 
to the II-line integral $\oint_\Sigma P(\xi,\eta)d\xi+Q(\xi,\eta)d\eta$, then using Green's formula to transform it to double integral $\iint_\Delta(\frac{\partial Q}{\partial \xi}-\frac{\partial P}{\partial \eta})d\xi d\eta$ on $\xi,\eta$-plane, following is how we 
calculate the double integral (often transform it to iterated integral $\int_{D\xi} d\xi \int_{\eta_0(\xi)}^{\eta(\xi)} (\frac{\partial Q}{\partial \xi}-\frac{\partial P}{\partial \eta})d\xi$, or others).\\
/\textcolor{skycyan}{Shortly,} calculate the area of $D$ on $x,y$ plane: coordinate transform with parameterization $\to$ Green's formula $\to$ Double integration calculate/ \\
\textcolor{skycyan}{Comparison}: \\
Recaping the previous $D=\frac 1 2 \oint_S x dy - y dx$, the original calculation method for this II-line integral is that find thier parametric equation $\begin{cases} x=x(t)\\ y=y(t) \end{cases}, t\in[a,b]$ for different (closed) curve $S$, then substitute into the 
calculation.From here, we could see that the essense of II-line integral is definite integral and has chirality.\\
However, "coordinate transform + Green's formula" method: $D=\iint_{\Delta}\frac{D(x,y)}{D(\xi,\eta)} d\xi d\eta$ as long as we master the coordinate transformation of a type of shape equation, we can easily calculate the results. Although there is a parameterization 
in the derivation, it is not reflected in the result (it can be seen from the process that if the parameterization is continued, it will be the result of the II-line integral in $\xi,\eta$-plane). In essence, the area is calculated as the volume with the same value, 
similiar with $\frac 1 2 \oint_S x dy-y dx=\iint_D 1 dxdy$ or $\oint_S x dy=\iint_D 1 dxdy$ (If you unfold the left and right sides under the curved trapezoid, you can find the consistency.), it's just that the physical meanings of the two equations are different.
\begin{figure}[H]
    \centering
    \includegraphics[scale=0.32]{"./Figures/250815_4.png"}
\end{figure}
/\textcolor{skycyan}{Note}: The choice of which method should be based on which is more convenient, parametric equations or coordinate transformations./
\end{minipage}}
\newline
\fbox{\begin{minipage}{0.95\textwidth}
\textcolor{skycyan}{Variable Substitution in Double Integral}\\
$$\iint_D f(x,y) dxdy=\begin{cases} x=x(\xi,\eta)\\ y=y(\xi,\eta) \end{cases}=\iint_{\Delta} f(x(\xi,\eta),y(\xi,\eta))|\frac{D(x,y)}{D(\xi,\eta)}| d\xi d\eta$$
\end{minipage}}
\subsection{Surface Integral}
\subsubsection{Definition and basics}
\fbox{\begin{minipage}{0.95\textwidth}
\textcolor{skycyan}{Definition}\\
$S=\lim\limits_{\lambda\to 0}\sum\limits_i T_i$, where $T_i$ is the projection area of $S_i$ on the tangent plane at point $M_i\in S_i$.
\begin{figure}[H]
    \centering
    \includegraphics[scale=0.36]{"./Figures/250815_5.png"}
\end{figure}
\end{minipage}}
\newline\newline
\fbox{\begin{minipage}{0.95\textwidth}
\textcolor{skycyan}{Calculte Surface Integral}(Omit matrix-vector operations)\\
$$S=\lim\limits_{\lambda\to 0}\sum\limits_i T_i=\iint_{\Delta} \sqrt{A^2+B^2+C^2} du dv=\iint_{\Delta} \sqrt{EG-F^2} du dv$$
where $A=\begin{vmatrix} y_u' & z_u'\\ y_v' & z_v' \end{vmatrix},\ B=\begin{vmatrix} z_u' & x_u'\\ z_v' & x_v' \end{vmatrix},\ C=\begin{vmatrix} x_u' & y_u'\\ x_v' & y_v' \end{vmatrix}$, 
$E=\sum\limits_{x,y,z} x_u'^2,\ F=\sum\limits_{x,y,z} x_u'x_v',\ G=\sum\limits_{x,y,z} x_v'^2$.
\end{minipage}}
\newline\newline
\fbox{\begin{minipage}{0.95\textwidth}
\textcolor{skycyan}{Procedure:} \\
$\forall M_1\in S'$, the coordinates of its projection $M_1'$ on $T'$ are $\begin{cases} \xi=\xi(u,v)\\ \eta=\eta(u,v)\\ \zeta=\zeta(u,v) \end{cases}$, 
then $S_T'=\iint_{\Delta}|\frac{D(\xi,\eta)}{D(u,v)}| du dv$, then the results could appear after a series of matrix-vector operations.
\begin{figure}[H]
    \centering
    \includegraphics[scale=0.36]{"./Figures/250815_6.png"}
\end{figure}
\end{minipage}}

\newpage
\section{Multiple Integrals}
\subsection{The Riemann Integral over an $n$-Dimensional Interval}
\subsubsection{Comprehensions}
First, the definition of Riemann integral over $n$-dimention interval is simimiar (almost same) with $1$-dimension condition ($\int_{I}f(x)dx := \lim\limits_{\lambda(P)\to 0}\sigma(f,P,\xi)$),
$$\int_{I}f(x^1,\cdots,x^n)dx^1\cdots dx^n\ \text{or}\ \int\cdots\int_{I}f(x^1,\cdots,x^2)dx^1\cdots dx^n$$
Then, introduce the necessity condition of it ($f\in\mathcal{R}(I)\implies f\ \text{is bounded on}\ I$).
After that, book gives some properties (Lemmas) about the set of measure zero (same with $1$-dimension condition), and developes the Cantor Theorem (about uniform continuity of continuous functions on compact set).
At last, we're given 2 criterions (necessity and sufficiency) of the Riemann integral (Darboux Criterion, Lebesgue Criterion). 
\newline\newline
\fbox{\begin{minipage}{0.95\textwidth}
\begin{ex*}
Let $f:I\to\mathbb{R}$ is a continuous real-value function on $I\subset\mathbb{R}^{n-1}$, then it's graph is a set measure zero in $n$-dimension.
\end{ex*}
\end{minipage}}
\newline
\fbox{\begin{minipage}{0.95\textwidth}
\begin{lem}
/Development of Cantor Theorem/ If the relation $\omega(f;x) < w_0$ holds at each point of a compact set 
$K$ for the function $f : K \to \mathbb{R}$, then for every $\varepsilon > 0$ there exists $\delta > 0$ such 
that $\omega(f; U_{\delta}(x)) < w_0 + \varepsilon$ for each point $x \in K$.
\end{lem}
\end{minipage}}
\newline
\fbox{\begin{minipage}{0.95\textwidth}
\begin{thm}
(Lebesgue's criterion). $f \in \mathcal{R}(I) \Leftrightarrow (f$ is bounded on $I) \land$ 
$(f$ is continuous almost everywhere on $I)$.
\end{thm}
\end{minipage}}
\newline
\fbox{\begin{minipage}{0.95\textwidth}
\begin{thm}
(Darboux Criterion) $f\in\mathcal(I)\Leftrightarrow(f\ is\ bounded\ on\ I)\wedge(\underline{\mathcal{J}}=\overline{\mathcal{J}})$, 
where $\underline{\mathcal{J}}=\sup\limits_P s(f,P),\ \overline{\mathcal{J}}=\inf\limits_P S(f,P)$.
\end{thm}
\end{minipage}}
\newpage
\subsubsection{Problems Sets}
\begin{prb}
\begin{enumerate}
  \item[a)] Show that a set of measure zero has no interior points.
  \item[b)] Show that not having interior points by no means guarantees that a set is of measure zero.
  \item[c)] Construct a set having measure zero whose closure is the entire space $\mathbb{R}^n$.
  \item[d)] A set $E \subset I$ is said to have content zero if for every $\varepsilon > 0$ it can be covered by a finite system of intervals $I_1, \ldots, I_k$ such that $\sum_{i=1}^k |I_i| < \varepsilon$. Is every bounded set of measure zero a set of content zero?
  \item[e)] Show that if a set $E \subset \mathbb{R}^n$ is the direct product $\mathbb{R} \times e$ of the line $\mathbb{R}$ and a set $e \subset \mathbb{R}^{n-1}$ of $(n-1)$-dimensional measure zero, then $E$ is a set of $n$-dimensional measure zero.
\end{enumerate}
\end{prb}
\fbox{\begin{minipage}{0.95\textwidth}
\begin{pf*}
\begin{enumerate}
  \item[a)] [Proof by contradiction] If $E$ has a point $x_0$ as an interior point, then there exists a ball $B(x_0, r)$ such that $B(x_0, r) \subset E$. Since the ball has positive measure, it follows that $E$ cannot have measure zero.
  \item[b)] [Counterexample] Consider the irrational numbers $\mathbb{R} \setminus \mathbb{Q}$, which have no interior points in $\mathbb{R}$ but are not of measure zero.
  \item[c)] /inspired by b)/ Consider the set $E = \mathbb{Q}^n$, the set of all points in $\mathbb{R}^n$ with rational coordinates. The closure of $E$ is $\mathbb{R}^n$, and $E$ has measure zero since it can be covered by countably many intervals whose total length can be made arbitrarily small.
  \item[d)] [Principle] No, since the Lebesgue measure zero set can be covered by at most countable intervals (which means measure $\mu=o(1)\cdot\mathcal{O}(1)$), while a set of content zero can be covered by a finite number of intervals (which means the content (also a kind of measure) sometimes is $\mathcal{O}(1)\cdot\mathcal{O}(1)$). [Example] The $\mathbb{Q}\cap[0,1]$ and Cantor set are sets of measure zero but none of content zero.
  \item[e)] From the construction of the product measure $\lambda_n$ (Lebesgue measure of $\mathbb{R}^n$) /Lemma1/ if $A \in B(\mathbb{R}^k)$ and $B \in B(\mathbb{R}^{n-k})$, ($0 < k < n$), then $\lambda_n(A \times B) = \lambda_k(A) \lambda_{n-k}(B)$. So here (using also that /Lemma2/ if $E_m \uparrow E$, $\lim_{m \to +\infty} \lambda_n(E_m) = \lambda_n(E)$) $\lambda_n(E) = \lambda_n(\mathbb{R} \times e)$ $= \lim_{m \to +\infty} \lambda_n((-m, m) \times e)$ 
  $= \lim_{m \to +\infty} \lambda_1((-m, m)) \times \lambda_{n-1}(e)$ $= \lim_{m \to +\infty} \lambda_1((-m, m)) \times 0$ $= 0$.
\end{enumerate}
\begin{nt*}
1. The proof of e) /Lemma2/ could found in \textit{Real and Complex Analysis} by Walter Rudin, Theorem 1.19(Page 16), based on the (positive) measure definition (countable additive of disjoint countable measurable sets).
\end{nt*}
\end{pf*}
\end{minipage}}
\newpage
\begin{prb}
\begin{enumerate}
  \item [a)] Construct the analogue of the Dirichlet function in $\mathbb{R}^{n}$ and show that a bounded function $f:I\to\mathbb{R}$ equal to zero at almost every point of the interval $I$ may still fail to belong to $\mathcal{R}(I)$.
  \item [b)] Show that if $f\in\mathcal{R}(I)$ and $f(x)=0$, at almost all points of the interval $I$, then $\int\limits_{I}f(x)\,{\rm d}x=0$.
\end{enumerate}
\end{prb}
\fbox{\begin{minipage}{0.95\textwidth}
\begin{pf*}
\begin{enumerate}
  \item[a)] [Counterexample] Consider $\mathcal{D}^n=\begin{cases} 1 & \text{ if } x\in\mathbb{Q}^n \\ 0 & \text{ if } x\in\mathbb{R}^n\backslash\mathbb{Q}^n \\\end{cases}$, the discontinuous set is $\mathcal{R}^n$.
  \item[b)] /by definition/ $\int_{I}f(x)dx=\lim\limits_{\lambda(P)\to 0}\sum\limits_{\xi_i\in P}f(\xi_i)\Delta(x_i)=\lim\limits_{\lambda(P)\to 0}\sum\limits_{\xi_i\in A\cap P}f(\xi_i)\Delta(x_i)+\lim\limits_{\lambda(P)\to 0}\sum\limits_{\xi_j\in B\cap P}f(\xi_j)\Delta(x_i)=\mathcal{O}(1)\cdot o(1)+o(1)\cdot\mathcal{O}(1)$, where $A=\{\cup[x_k,x_{k+1}]\subset P:\forall k,\ \exists \xi\in[x_k,x_{k+1}],\ s.t.\ f(\xi)\neq 0\}$ and $B=I\backslash A$. 
\end{enumerate}
\end{pf*}
\end{minipage}}
\begin{prb}
There is a small difference between our earlier definition of the Riemann integral on a closed interval $I \subset \mathbb{R}$ and Definition 7 for the integral over an interval of arbitrary dimension. This difference involves the definition of a partition and the measure of an interval of the partition. Clarify this nuance for yourself and verify that
$$\int_a^b f(x) dx = \begin{cases} \int_I f(x) dx & \text{ if } a < b \\ - \int_I f(x) dx & \text{ if } a > b \end{cases}$$
where $I$ is the interval on the real line $\mathbb{R}$ with endpoints $a$ and $b$.
\end{prb}
\fbox{\begin{minipage}{0.95\textwidth}
\begin{pf*}
Just understand it, it's not hard.
%The difference is that in the earlier definition, the partition is a finite set of points dividing the interval into subintervals, while in Definition 7, the partition can be any finite collection of intervals. The measure of an interval in $\mathbb{R}$ is simply its length, which is the difference between its endpoints.
%To verify the integral equality, we can consider the definition of the Riemann integral over an interval $I = [a, b]$:
%$$\int_I f(x) dx = \lim_{\lambda(P) \to 0}\sum_{\xi_i \in P} f(\xi_i) \Delta(x_i)$$
%where $P$ is a partition of the interval $I$ and $\Delta(x_i)$ is the length of the subintervals. If $a < b$, then the integral is computed directly over the interval $[a, b]$. If $a > b$, we can redefine the integral as:
%$$\int_a^b f(x) dx = -\int_b^a f(x) dx = -\lim_{\lambda(P) \to 0}\sum_{\xi_i \in P} f(\xi_i) \Delta(x_i)$$
%where $P$ is a partition of the interval $[b, a]$. This shows that the integral is negative when the limits are reversed, confirming the equality.
\end{pf*}
\end{minipage}}
\begin{prb}
\begin{enumerate}
  \item [a)] Prove that a real-valued function $f : I \to \mathbb{R}$ defined on an interval $I \subset \mathbb{R}^n$ is integrable over that interval if and only if for every $\varepsilon > 0$ there exists a partition $P$ of $I$ such that $S(f; P) - s(f; P) < \varepsilon$.
  \item [b)] Using the result of a) and assuming that we are dealing with a real-valued function $f : I \to \mathbb{R}$, one can simplify slightly the proof of the sufficiency of the Lebesgue criterion. Try to carry out this simplification by yourself.
\end{enumerate}
\end{prb}
\fbox{\begin{minipage}{0.95\textwidth}
\begin{pf*}
\begin{enumerate}
  \item[a)] /by The Darboux Criterion/ $f\in\mathcal(I)\Leftrightarrow(f\ is\ bounded\ on\ I)\wedge(\underline{\mathcal{J}}=\overline{\mathcal{J}})$, where $\underline{\mathcal{J}}=\sup\limits_P s(f,P),\ \overline{\mathcal{J}}=\inf\limits_P S(f,P)$. Notice that $\forall \varepsilon >0,\ \exists P,\ s.t.\ S(f,P)-s(f,P)<\varepsilon \Leftrightarrow \underline{\mathcal{J}}=\overline{\mathcal{J}}$ (proof is similiar with 1-dimention condition), which is what we need to prove.
  \item[b)] [Segmented estimation] Sufficiency: procedure is same with book, a) just simplify the last step.
\end{enumerate}
\end{pf*}
\end{minipage}}


\end{document}